\chapter{Pengantar}

Apakah ini salah prabowo kita menjadi miskin? Hidup miskin bukan berarti kita terbebas dari tangggung jawab untuk mendesain IC tapi dengan tingkat kemisikinan tertentu, mana bisa membeli software harga puluhan ribu dollar. Proffesor UI ajah belum tentu mampu apa lagi kalian para masyarakat dhuafa.

Sekali lagi mendesain IC itu bukan sekadar berapa tebal dompet anda, buktinya Bahlil yang dompetnya tebal saja tidak bisa mendesain IC. 

Kita tahu bahwa pemerintah indonesia membeli lisensi arm dengan harga jutaan dollar. Padahal kita bisa bikin processor sendiri, saya berfikir positif saja memang karena Indonesia memang negara yang kaya. Buku ini akan memberikan gambaran bagi kaum dhuafa jika ingin menjual sebuah arsitektur ke Indonesia sebaiknya mengaku sebagai orang Amerika. Kenapa? Karena orang indonesia itu Xenocentrism, suka ngomong antek asing tapi sebenernya intel Amerika, jelas maksud saya bukan Prabowo. Tapi diri kita yang suka dengan produk Amerika, kita adalah agen-agen Iphone, agen-agen tesla, bahkan mungkin kita adalah agen arm. Supaya apa? Terlihat lebih keren. Indonesia yang selalu ingin tampil branded secara singkat lebih suka ngikutin apa hegemoni ketimbang effort untuk bikin branding sendiri, namanya juga Xenosentrisme, biar keren, biar dihormati, biar didengar dan biar tidak diremehkan. Tapi itu hanya pinjaman, pinjaman branding dari amerika. Pinjaman yang suatu saat harus dikembalikan. Buat saya, saya tidak mempermasalahkan kita ingin terlihat keren, tapi terlihat keren dengan memimjam branding asing itu benar benar katrok.  